% This file is embedded in datatool-user.pdf version 3.0.1 2025-03-05
% Example 69 Automatically Formatting Values Calculated by Actions
% Label: "ex:actiondatum"
% arara: pdflatex
% arara: pdfcrop
\documentclass[12pt]{article}
\pagestyle{empty}
\usepackage{datatool}
% custom expandable command:
\newcommand{\limiteded}{limited edition}
% define a database with the name 'pricelist':
\DTLsetup{store-datum,default-name=pricelist}
\DTLaction{new}% create the default database
% 1st row:
\DTLaction[
 assign={
  Product = {The Adventures of Duck and Goose},
  Quantity = {1,452}, Price = {\$1.99}
 }
]{new row}
% 2nd row:
\DTLaction[
 assign={
  Product = {Duck and Goose on Holiday},
  Quantity = {94}, Price = {\$2.99}
 }
]{new row}
% the next value needs to be expanded:
\DTLaction[
key={Notes}, expand-value={\limiteded} ]{new entry}
% 3rd row:
\DTLaction[
 assign={
  Product = {The Return of Sir Quackalot},
  Quantity = {3}, Price = {\$4.99}
 }
]{new row} 
\begin{document}
 \section{Default datum=false} 

\DTLaction[
 key=Quantity,
 key2=Price,
 options={sd,min,max}
]{aggregate} 

Quantity column index: \DTLuse{column}.
Total quantity: \DTLuse{sum}.
Average quantity: \DTLuse{mean}.
Quantity standard deviation: \DTLuse{sd}.
Minimum quantity: \DTLuse{min}.
Maximum quantity: \DTLuse{max}. 

Price column index: \DTLuse{column2}.
Total price: \DTLuse{sum2}.
Average price: \DTLuse{mean2}.
Price standard deviation: \DTLuse{sd2}.
Minimum price: \DTLuse{min2}.
Maximum price: \DTLuse{max2}. 

\section{datum=\{round=2\}} 

\DTLaction[
 datum={round=2},
 key=Quantity,
 key2=Price,
 options={sd,min,max}
]{aggregate} 

Quantity column index: \DTLuse{column}.
Total quantity: \DTLuse{sum}
(\DTLget[sum]{\theTotal}\DTLdatumvalue{\theTotal}).
Average quantity: \DTLuse{mean}.
Quantity standard deviation: \DTLuse{sd}.
Minimum quantity: \DTLuse{min}
(\DTLget[min]{\theMin}\DTLdatumvalue{\theMin}).
Maximum quantity: \DTLuse{max}
(\DTLget[max]{\theMax}\DTLdatumvalue{\theMax}). 

Price column index: \DTLuse{column2}.
Total price: \DTLuse{sum2}.
Average price: \DTLuse{mean2}.
Price standard deviation: \DTLuse{sd2}.
Minimum price: \DTLuse{min2}.
Maximum price: \DTLuse{max2}. 

\section{datum=\{round=2,currency\}} (All aggregate values become currency!) 

\DTLaction[
 datum={round=2,currency},
 key=Quantity,
 key2=Price,
 options={sd,min,max}
]{aggregate} 

Quantity column index: \DTLuse{column}.
Total quantity: \DTLuse{sum}
(\DTLget[sum]{\theTotal}\DTLdatumvalue{\theTotal}).
Average quantity: \DTLuse{mean}.
Quantity standard deviation: \DTLuse{sd}.
Minimum quantity: \DTLuse{min}
(\DTLget[min]{\theMin}\DTLdatumvalue{\theMin}).
Maximum quantity: \DTLuse{max}
(\DTLget[max]{\theMax}\DTLdatumvalue{\theMax}). 

Price column index: \DTLuse{column2}.
Total price: \DTLuse{sum2}.
Average price: \DTLuse{mean2}.
Price standard deviation: \DTLuse{sd2}.
Minimum price: \DTLuse{min2}.
Maximum price: \DTLuse{max2}. 
\end{document}
